% 2016 年全国卷Ⅱ文科第 7 题：圆柱与圆锥组合几何体的三视图
\begin{tikzpicture}[>=Stealth, scale=0.9, line width=0.45pt, font=\small]
  % 两幅立面图：圆柱矩形与圆锥三角形共用底边。
  \draw (0,0) rectangle (2,2);
  \draw (0,2) -- (1,{2+sqrt(3)}) -- (2,2);
  \draw (3,0) rectangle (5,2);
  \draw (3,2) -- (4,{2+sqrt(3)}) -- (5,2);

  % 第一幅立面左侧的高度尺寸。
  \draw (-0.06,0) -- (-0.38,0);
  \draw (-0.06,2) -- (-0.38,2);
  \draw[<->] (-0.38,0) -- (-0.38,2);
  \node[left=2pt] at (-0.38,1) {$4$};
  \draw (-0.06,{2+sqrt(3)}) -- (-0.38,{2+sqrt(3)});
  \draw[<->] (-0.38,2) -- (-0.38,{2+sqrt(3)});
  \node[left=2pt] at (-0.38,{2+sqrt(3)/2}) {$2\sqrt{3}$};

  % 两幅立面图下方的宽度尺寸。
  \draw (0,-0.06) -- (0,-0.36);
  \draw (2,-0.06) -- (2,-0.36);
  \draw[<->] (0,-0.36) -- (2,-0.36);
  \node[below=1pt] at (1,-0.36) {$4$};
  \draw (3,-0.06) -- (3,-0.36);
  \draw (5,-0.06) -- (5,-0.36);
  \draw[<->] (3,-0.36) -- (5,-0.36);
  \node[below=1pt] at (4,-0.36) {$4$};

  % 俯视图。
  \draw (1,-2.15) circle (1);
  \fill (1,-2.15) circle (1.2pt);
\end{tikzpicture}
